\documentclass[onecolumn, draftclsnofoot,10pt, compsoc]{IEEEtran}
\usepackage{graphicx}
\usepackage{url}
\usepackage{setspace}
\usepackage{listings}

\usepackage{geometry}
\geometry{textheight=9.5in, textwidth=7in}

% 1. Fill in these details
\def \CapstoneTeamName{Manovacumeter Team}
\def \CapstoneTeamNumber{20}
\def \GroupMemberOne{Cade Raichart}
\def \GroupMemberTwo{Daniel Kato}
\def \GroupMemberThree{Nathan Shepherd}
\def \CapstoneProjectName{ISS Barometer App }
\def \CapstoneSponsorCompany{NASA}
\def \CapstoneSponsorPerson{Don Pettit}

% 2. Uncomment the appropriate line below so that the document type works
\def \DocType{Progress Report}

\newcommand{\NameSigPair}[1]{\par
\makebox[2.75in][r]{#1} \hfil 	\makebox[3.25in]{\makebox[2.25in]{\hrulefill} \hfill		\makebox[.75in]{\hrulefill}}
\par\vspace{-12pt} \textit{\tiny\noindent
\makebox[2.75in]{} \hfil		\makebox[3.25in]{\makebox[2.25in][r]{Signature} \hfill	\makebox[.75in][r]{Date}}}}
% 3. If the document is not to be signed, uncomment the RENEWcommand below
%\renewcommand{\NameSigPair}[1]{#1}

%%%%%%%%%%%%%%%%%%%%%%%%%%%%%%%%%%%%%%%
\begin{document}
\begin{titlepage}
    \pagenumbering{gobble}
    \begin{singlespace}
    	\includegraphics[height=4cm]{coe_v_spot1}
        \hfill
        % 4. If you have a logo, use this includegraphics command to put it on the coversheet.
        \includegraphics[height=4cm]{NASA}
        \par\vspace{.2in}
        \centering
        \scshape{
            \huge CS Capstone \DocType \par
            {\large\today}\par
            \vspace{.5in}
            \textbf{\Huge\CapstoneProjectName}\par
            \vfill
            {\large Prepared for}\par
            \Huge \CapstoneSponsorCompany\par
            \vspace{5pt}
            {\Large\NameSigPair{\CapstoneSponsorPerson}\par}
            {\large Prepared by }\par
            Group\CapstoneTeamNumber\par
            % 5. comment out the line below this one if you do not wish to name your team
            %\CapstoneTeamName\par
            \vspace{5pt}
            {\Large
                \NameSigPair{\GroupMemberOne}\par
                \NameSigPair{\GroupMemberTwo}\par
                \NameSigPair{\GroupMemberThree}\par
            }
            \vspace{20pt}
        }
        \begin{abstract}
        % 6. Fill in your abstract
        The International Space Station, the conglomerate work of multiple national space associations, including NASA, requires precision devices to monitor the many aspects of space life.
        One very important aspect for the success of life in near vacuum is sustained air pressure inside the modules.
        This air pressure is monitored by mechanical barometers, called Manovacometers, which must be visually monitored, and are no longer available.
        The \CapstoneProjectName seeks to provide a suitable replacement for the current barometer device aboard the International Space Station.
        This document will serve to introduce the \CapstoneProjectName, and will detail the problem, the solution, and the success metrics of the project.
        The sections include the Definition of the Problem, Definition of Project, and the Metrics of Success.
        \end{abstract}
    \end{singlespace}
\end{titlepage}
\newpage
\pagenumbering{arabic}
\tableofcontents
% 7. uncomment this (if applicable). Consider adding a page break.
%\listoffigures
%\listoftables
\clearpage

% 8. now you write!
\section{Project Purpose}

\section{Project Goals}

\section{Current State of the Project}
% If anyone can think of a better way to state this, feel free to change the heading

\section{Project Problems}
We chose to use swift and xcode because they allows for quick development and prototyping. 
This said we ran into many problems with that. For one Swift only works on Linux and MacOS.
One of our group members did not have either at the time. They fixed this with installing a linux OS. 
Another problem was installing the swift libraries and connecting them with the xcode application.
Linux needed many packages to be downloaded for the library to work correctly, this took many hours of troubleshooting as it was not straightforward.
Connecting swift to xcode on the MacOS was another problem as well. 
Our last problem with swift and xcode is we have no experience using either.
We had and still have to continually work on improving our knowledge of these systems using online videos, guides and other resources. 

For the first time we all experienced what it was like working with someone who is in the workplace.
We learned that our clients will not be getting back to us immediately and sometimes we have to wait for them reply.
This means we have to work and finish our work and send it to him way earlier then we had planned originally. 
We also ran into the problems of finding the correct IEEE documentations. 
Sometimes we found the document quick and easy while other times it hours and many students combined to find them.  

\section{Technology Review Paper}

The goal of this document is to justify the choices made, with research to prove that our project was done in the best way possible. 
Included is a list of technology options, details on each, and justifications of each choice. 
The nine technologies we needed to find implementations  for was development platform, chart library, UI organization, barometer implementation, data storage, graph library, color palette, programing language and UI library. 
We will summaries the appropriate implementations that we chose and the justification for those implementations. 

For our development platform technology we will be using Xcode and swift to develop the ISS Barometer application.
We chose this technology for its ability to interface with the onboard barometer, which is required for this project.
The ability to interface directly with the built in barometer is necessary for gaining an accurate pressure reading.[2 D


Our Chart Library Technology will be the library called Charts.
Reliability is of the utmost importance for our application. 
Charts offers support for dynamic graphs, while the other 2 libraries don't have the consistent support that will be need.
The Charts API also makes it easy to add points to a dynamic graph, making it the perfect library. [2 D]

We are using the single page layout for our Organization of our UI. 
Using a single page layout will allow us to access our settings, statistics and our graph frequently and easily.
Since our astronauts will be using this in emergency situations speed and easy access is absolutely necessary.  [2 D]

Our client has asked us to us the iPad air 2 for our Barometer Implementation technology.
We are asked to use this because the space station has 8-10 iPads aboard. 
Using any other implementation would cause us to ship it up to the space station which is not necessary or convenient. [3 C]

For the Graph Interface technology we have chosen Swift Chars library.
Our client has requested that the graph be able to zoom and pan, the only library that allows us to do this is the swift charts library.  [3 c]

The UI Library will be mostly Cocoa controls for is massive library and larger features list. 
Though as we have found no single library has everything we need to complete our task, therefore we will need to use many other libraries such as CRPageViewController and FlatUIKit. [3 c]
    

We will be programing in the swift Language. 
This was chosen for its ability to quickly prototype and develop the application. 
Swift also has many performance benefits for app creation that make it the best option. 
Also our client doesn't know much about programming swift works xcode which will make showing him our progress easily. [1 nathan]

For our data storage technology we will use CoreData, in tandem with NSUserDefaults.
CoreData allows for large data sets that scale well and far better than the other libraries do.  
Using NSUserDefaults package, allows us to quickly save data, such as float point accuracy. [1 nathan]

Our color palette technology will be the Subtle Technology Theme.
This was chosen because the other options, OSU similar theme and Nasa similar theme, are designed for specific uses that would be intrusive and counterproductive in our application.
Subtle Technology Theme allows the user to not be distracted by unwanted color.  [1 nathan]

\section{Interesting Code}
\begin{lstlisting}
  func startDisplayingPressureData() {
      self.altimeter.startRelativeAltitudeUpdates(to: OperationQueue.main, withHandler: {
        data, error in
          self.handlePressureReading(data: data!)
      })
  }

\end{lstlisting}

One of the main features of our application is its ability to read data from the iPad's built in barometer and display the data collected in different ways.
The way that data is collected from the barometer is quite interesting.
The process of reading data from the barometer requires a bit of processing power and time which can block other time sensitive tasks form being performed.
As a result, the measurement is done asynchronously through an object called the \texttt{OperationQueue}.
A method provided by the CoreMotion library called \texttt{startRelativeUpdates()} takes the \texttt{OperationQueue}'s \texttt{main} queue, and a callback function as arguement.
The function then appends a task to the queue to perform in the background, and when the task is complete, the result is passed into the callback function for further processing.\\

\section{Retrospective}
\begin{center}
	\begin{tabular}{| l | p{0.3\linewidth} | p{0.3\linewidth} | p{0.3\linewidth} |}
		\hline
		Week 1 & Positives & Deltas & Actions \\ \hline
		Week 2 & & & \\ \hline				% Put content between '&'s
		Week 3 & Finished problem statements and turned into both our client and teacher
 & Figure out how to share app with client, whom does not know the github system
 & step by step instructions on how to download the app \\ \hline
		Week 4 & & & \\ \hline
		Week 5 & & & \\ \hline
		Week 6 &First prototype with two basic display of barometer numbers
  finished the Requirements Document & & \\ \hline
		Week 7 & & & \\ \hline
		Week 8 & Completed Technology Review document rough draft & Need to complete Technology Review rough draft & A rough draft of the Technology Review will be created \\ \hline
		Week 9 &Finished Technology document
Started working on design document
 & Could not find IEEE paper & Searched more and asked around\\ \hline
		Week 10 & & & \\ \hline




	\end{tabular}
\end{center}




\end{document}
